% SOURCE_STATEMENT: "For every asymptotically flat Riemannian three-manifold $(M^3,g)$, does the capacity–volume mass equal the ADM mass, $m_{\\mathrm{CV}}(M,g)=m_{\\mathrm{ADM}}(M,g)$?"
% Preparation and final self-audit: CONTRIBUTING.md prompts C and D; source/audit artifacts retained in the local physics_statements_2026_10_11 batch.
% !TeX program = lualatex
\documentclass[11pt]{article}
\usepackage[margin=25mm]{geometry}
\usepackage{fontspec}
\setmainfont{Latin Modern Roman}
\usepackage{amsmath,amssymb,mathtools}
\usepackage{unicode-math}
\setmathfont{Latin Modern Math}
\usepackage{xurl,hyperref}
\hypersetup{colorlinks=true,urlcolor=blue}
\setcounter{secnumdepth}{0}
\setlength{\parindent}{0pt}
\setlength{\parskip}{5pt}
\setlength{\emergencystretch}{3em}
\newcommand{\sref}[2]{\hyperlink{source:#1:#2}{[S#2]}}
\newcommand{\cataloguescope}[1]{\paragraph{Recorded target.}#1}
\newcommand{\R}{\mathbb R}
\newcommand{\C}{\mathbb C}
\newcommand{\N}{\mathbb N}
\newcommand{\Z}{\mathbb Z}
\newcommand{\Q}{\mathbb Q}
\newcommand{\F}{\mathbb F}
\newcommand{\eps}{\varepsilon}
\begin{document}
\section{Equality of Capacity–Volume and ADM Mass}\label{statement}

\subsection{Definitions and mathematical statement}

\paragraph{Shared notation.}$\mathbb N=\{1,2,\ldots\}$; the standard number fields have their usual meanings. All additional conventions are stated locally.


An asymptotically flat three-manifold approaches Euclidean space on its end. The ADM mass is the asymptotic flux of first derivatives of the metric. The normalized capacity of a compact region $K$ is the infimum of $(4\pi)^{-1}\int_{M\setminus K}|\nabla u|^2dV_g$ over admissible functions with boundary value one on $K$ and value zero at infinity. Writing $|K|$ for Riemannian volume, Jauregui's capacity--volume mass is
\[
m_{\mathrm{CV}}(M,g)=\sup_{\{K_j\}}\limsup_{j\to\infty}
\left[\left(\frac{3|K_j|}{4\pi}\right)^{1/3}-\operatorname{cap}(K_j)\right],
\]
where the supremum is over the admissible compact exhaustions in the source.

\paragraph{Catalogue wording retained for traceability.}
For every asymptotically flat Riemannian three-manifold $(M^3,g)$, does the capacity–volume mass equal the ADM mass, $m_{\mathrm{CV}}(M,g)=m_{\mathrm{ADM}}(M,g)$?

\paragraph{Scope and source qualification.}The source reports equality for harmonically flat ends under additional hypotheses, and asks about general asymptotically flat manifolds. An exact proof must use the report's definition of the exhaustion supremum, rather than a single chosen family of regions.
\subsection{Short English statement}
Determine whether the mass defined from large-region capacity and volume always equals ADM mass.
\subsection{Sources}
\begin{itemize}
\item[\hypertarget{source:statement:1}{S1}] ULAM.AI and the UnsolvedMath Contributors, UnsolvedMath, catalogue entry. Catalogue curation under CC BY 4.0; definitions and English exposition added locally.
\url{https://huggingface.co/datasets/ulamai/UnsolvedMath}.
\item[\hypertarget{source:statement:2}{S2}] Mathematical Aspects of General Relativity (2021), Oberwolfach Reports, DOI 10.4171/owr/2021/40, pp. 2157--2267. Original report and mathematical context.
\url{https://doi.org/10.4171/owr/2021/40}.
\end{itemize}
\label{end:statement}
\end{document}
